\documentclass[9pt,a4paper,twocolumn]{ctexart}

% --- Packages ---
\usepackage[top=1.2cm, bottom=1.2cm, left=1.2cm, right=1.2cm, columnsep=0.8cm]{geometry}
\usepackage{hyperref}
\usepackage{graphicx}
\usepackage{float}
\usepackage{longtable}
\usepackage{tabularx}
\usepackage{booktabs}
\usepackage{enumitem}
\usepackage{xcolor}
\usepackage{fancyhdr}
\usepackage{listings}
\usepackage{fontspec}
\usepackage{titlesec}
\usepackage{caption}

% --- Code block styling ---
\definecolor{codebg}{RGB}{245,245,245}
\definecolor{codeframe}{RGB}{220,220,220}
\definecolor{codekw}{RGB}{0,0,192}
\definecolor{codecomment}{RGB}{0,128,0}
\definecolor{codestring}{RGB}{163,21,21}
\definecolor{codenum}{RGB}{128,0,128}
\definecolor{codepunct}{RGB}{90,90,90}
\definecolor{badgegreen}{RGB}{34,139,34}
\definecolor{badgeblue}{RGB}{30,144,255}
\definecolor{badgeblack}{RGB}{50,50,50}

% --- Difficulty badges (rounded corners via tikz, pre-rendered in saveboxes) ---
\usepackage{tikz}
\newsavebox{\badgechubox}
\savebox{\badgechubox}{\tikz[baseline]{\node[fill=badgegreen,text=white,rounded corners=2.5pt,inner xsep=4pt,inner ysep=1.5pt,font=\sffamily\footnotesize\bfseries]{初级};}}
\newsavebox{\badgezhongbox}
\savebox{\badgezhongbox}{\tikz[baseline]{\node[fill=badgeblue,text=white,rounded corners=2.5pt,inner xsep=4pt,inner ysep=1.5pt,font=\sffamily\footnotesize\bfseries]{中级};}}
\newsavebox{\badgegaobox}
\savebox{\badgegaobox}{\tikz[baseline]{\node[fill=badgeblack,text=white,rounded corners=2.5pt,inner xsep=4pt,inner ysep=1.5pt,font=\sffamily\footnotesize\bfseries]{高级};}}
\newcommand{\lvchu}{\usebox{\badgechubox}}
\newcommand{\lvzhong}{\usebox{\badgezhongbox}}
\newcommand{\lvgao}{\usebox{\badgegaobox}}
\pdfstringdefDisableCommands{%
  \def\lvchu{[初级]}%
  \def\lvzhong{[中级]}%
  \def\lvgao{[高级]}%
}

\lstset{
  basicstyle=\ttfamily\scriptsize,
  keywordstyle=\color{codekw}\bfseries,
  commentstyle=\color{codecomment}\itshape,
  stringstyle=\color{codestring},
  backgroundcolor=\color{codebg},
  frame=single,
  rulecolor=\color{codeframe},
  breaklines=true,
  breakatwhitespace=false,
  breakindent=0pt,
  postbreak=\mbox{\textcolor{red}{$\hookrightarrow$}\space},
  numbers=left,
  numberstyle=\tiny\color{gray},
  columns=flexible,
  keepspaces=true,
  showstringspaces=false,
  tabsize=2,
  aboveskip=4pt,
  belowskip=4pt,
  extendedchars=true,
  literate={，}{{，}}1 {；}{{；}}1 {！}{{！}}1 {？}{{？}}1 {“}{{“}}1 {”}{{”}}1 {（}{{（}}1 {）}{{）}}1 {《}{{《}}1 {》}{{》}}1,
}

% --- Extra language definitions (no built-in listings support) ---
\lstdefinelanguage{json}{
  morekeywords={true,false,null},
  sensitive=true,
  morestring=[b]",
  comment=[l]{//},
  morecomment=[s]{/*}{*/},
}
\lstdefinelanguage{yaml}{
  morekeywords={true,false,null,yes,no,on,off},
  sensitive=false,
  morecomment=[l]{\#},
  morestring=[b]",
  morestring=[d]',
}
\lstdefinelanguage{http}{
  morekeywords={GET,POST,PUT,DELETE,PATCH,HEAD,OPTIONS,HTTP,Host,Accept,Authorization,Content-Type,Content-Length,User-Agent},
  sensitive=true,
  morecomment=[l]{\#},
  morestring=[b]",
}
\lstdefinelanguage{TypeScript}{
  morekeywords={abstract,any,as,async,await,break,case,catch,class,const,continue,debugger,declare,default,delete,do,else,enum,export,extends,false,finally,for,from,function,get,if,implements,import,in,instanceof,interface,keyof,let,module,namespace,never,new,null,number,of,package,private,protected,public,readonly,return,set,static,string,super,switch,symbol,this,throw,true,try,type,typeof,undefined,unknown,var,void,while,with,yield,Record,Array,Promise,AsyncIterable,ReadableStream,Date,Error,Object,JSON},
  sensitive=true,
  morecomment=[l]{//},
  morecomment=[s]{/*}{*/},
  morestring=[b]",
  morestring=[b]',
  morestring=[b]`{`},
}

% --- Compact spacing ---
\linespread{0.95}
\setlength{\parskip}{1pt plus 1pt minus 1pt}
\setlength{\parindent}{0pt}

% --- Table styling ---
\renewcommand{\arraystretch}{1.1}

% --- Heading styling (numbered) ---
\titleformat{\section}{\normalsize\bfseries}{\thesection\quad}{0em}{}
\titlespacing{\section}{0pt}{6pt}{2pt}
\titleformat{\subsection}{\small\bfseries}{\thesubsection\quad}{0em}{}
\titlespacing{\subsection}{0pt}{4pt}{1pt}
\titleformat{\subsubsection}{\small\bfseries}{\thesubsubsection\quad}{0em}{}
\titlespacing{\subsubsection}{0pt}{3pt}{1pt}

% --- Page style ---
\pagestyle{fancy}
\fancyhf{}
\fancyhead[L]{\small\emph{\leftmark}}
\fancyhead[R]{\small\thepage}
\renewcommand{\headrulewidth}{0.4pt}
\renewcommand{\sectionmark}[1]{}  % prevent sections from overwriting leftmark

% --- Hyperref setup ---
\hypersetup{
  colorlinks=true,
  linkcolor=blue,
  urlcolor=blue,
  citecolor=blue,
}

% --- Caption format ---
\DeclareCaptionFormat{myformat}{\small\bfseries #1#2#3}
\captionsetup{format=myformat}

\begin{document}

\title{Agent 案例库与反模式索引}
\date{}
\maketitle
\markboth{Python 版 Agent 知识}{ Python 版 Agent 知识}

\begin{quote}
语言策略：shared（Java canonical，P/T 同文）
\end{quote}



\section*{编号规则}


案例编号规则：\texttt{CASE-领域-序号}；反模式编号规则：\texttt{AP-序号}。新增案例或反模式必须登记到本节。



\section*{1. 案例库索引}

{\footnotesize
\begin{tabular}{|p{\dimexpr 0.107\columnwidth-2\tabcolsep\relax}|p{\dimexpr 0.259\columnwidth-2\tabcolsep\relax}|p{\dimexpr 0.196\columnwidth-2\tabcolsep\relax}|p{\dimexpr 0.438\columnwidth-2\tabcolsep\relax}|}
\hline
\textbf{编号} & \textbf{案例} & \textbf{关键做法} & \textbf{详细位置} \\
\hline
CASE-ENG-01 & OpenAI 三人团队 & 小团队 + 成熟 Harness，自动化优先 & \href{../04-工程实践/03-Harness工程.md}{Harness Engineering} \\
\hline
CASE-ENG-02 & Anthropic 三智能体 & 从上下文焦虑转向三智能体协作架构 & \href{../04-工程实践/03-Harness工程.md}{Harness Engineering} \\
\hline
CASE-ENG-03 & Stripe 每周 1300+ PR & 混合状态机编排，无人值守 PR & \href{../04-工程实践/03-Harness工程.md}{Harness Engineering} \\
\hline
CASE-ENG-04 & Mitchell Hashimoto 单人 Harness & 个人级约束、工具和上下文管理 & \href{../04-工程实践/03-Harness工程.md}{Harness Engineering} \\
\hline
CASE-LOOP-01 & Claude Code 的 /loop 与 /goal & 目标驱动的循环与自动修复边界 & \href{../04-工程实践/02-Loop工程.md}{Loop Engineering} \\
\hline
CASE-MEM-01 & Claude Code 的 CLAUDE.md & Markdown 作为项目记忆 & \href{../02-Agent/02-Agent记忆系统.md}{Agent 记忆系统} \\
\hline
\end{tabular}
}



\section*{2. 可复用模式}

{\footnotesize
\begin{tabular}{|p{\dimexpr 0.302\columnwidth-2\tabcolsep\relax}|p{\dimexpr 0.488\columnwidth-2\tabcolsep\relax}|p{\dimexpr 0.209\columnwidth-2\tabcolsep\relax}|}
\hline
\textbf{模式} & \textbf{一句话描述} & \textbf{适用场景} \\
\hline
小团队 + Harness & 不扩人，先扩自动化与护栏 & 早期 AI 产品 \\
\hline
混合状态机编排 & 固定流程管骨架，LLM 管节点内判断 & 高风险流程 \\
\hline
三智能体协作 & 生成、审查、验证角色分离 & 代码生成 \\
\hline
Markdown 记忆 & 项目事实放在版本库中 & 单人/团队开发工具 \\
\hline
目标循环 & /goal 定目标，/loop 受约束执行 & 多步自动化 \\
\hline
\end{tabular}
}



\section*{3. 反模式索引}

{\footnotesize
\begin{tabular}{|p{\dimexpr 0.053\columnwidth-2\tabcolsep\relax}|p{\dimexpr 0.200\columnwidth-2\tabcolsep\relax}|p{\dimexpr 0.105\columnwidth-2\tabcolsep\relax}|p{\dimexpr 0.642\columnwidth-2\tabcolsep\relax}|}
\hline
\textbf{编号} & \textbf{反模式} & \textbf{后果} & \textbf{对策与来源} \\
\hline
AP-01 & 目标写得太虚 & Agent 反复绕圈 & 目标必须可验收，\href{../04-工程实践/02-Loop工程.md}{Loop Engineering} \\
\hline
AP-02 & 把 Agent 自评当验收 & 错误自动通过 & 结构化验收标准 + 独立评审，\href{../04-工程实践/02-Loop工程.md}{Loop Engineering} \\
\hline
AP-03 & 没有成本上限 & 一次任务烧光预算 & Token/轮次/金额硬上限，\href{../04-工程实践/02-Loop工程.md}{Loop Engineering} \\
\hline
AP-04 & 权限给得太大 & 高危动作误执行 & L0–L3 工具分级 + 审批，\href{../06-安全与治理/01-Agent安全护栏清单.md}{Agent 安全护栏清单} \\
\hline
AP-05 & 第一版就自动修 & 小错误变大事故 & 先输出建议，稳定后再自动执行，\href{../04-工程实践/02-Loop工程.md}{Loop Engineering} \\
\hline
AP-06 & 上下文喂得越多越好 & 40\\% 后效果下降 & 上下文预算与裁剪，\href{../04-工程实践/03-Harness工程.md}{Harness Engineering} \\
\hline
AP-07 & 没有降级路径 & 模型失败系统不可用 & Fallback 与固定 Workflow 兜底，\href{../04-工程实践/04-大模型网关.md}{大模型网关} \\
\hline
AP-08 & 为了多 Agent 而多 Agent & 通信与调试成本爆炸 & 先用单 Agent + Workflow，\href{../02-Agent/03-多智能体编排.md}{多智能体编排} \\
\hline
AP-09 & 改 Prompt 不跑回归 & 线上静默退化 & 离线黄金集 + 回归门禁，\href{../05-评测与质量/01-Agent评测体系.md}{Agent 评测体系} \\
\hline
AP-10 & 没有 Trace & 事故无法复盘 & 决策轨迹、工具调用、上下文留痕，\href{../04-工程实践/05-Agent可观测与追踪.md}{Agent 可观测与追踪} \\
\hline
AP-11 & 一次性全量发布 & 灰度问题影响全量用户 & 影子流量 → 灰度 → 全量，\href{../04-工程实践/06-Agent部署与发布.md}{Agent 部署与发布} \\
\hline
AP-12 & 工具 description 只写名字 & 模型乱调工具 & 写清何时用、何时不用，\href{../02-Agent/01-Agent核心概念.md}{Agent 核心概念} \\
\hline
\end{tabular}
}



\section*{4. 新增案例模板}


\begin{lstlisting}
## CASE-XXX-01 案例名
- 背景：一句话说明业务场景
- 架构：使用的编排模式与工具
- 效果：可量化指标
- 经验：可复用的 1-3 条
- 反模式：踩过的坑（登记到 AP）
- 来源：文档或链接
\end{lstlisting}



\section*{5. 延伸阅读}

\begin{itemize}[itemsep=1pt, topsep=2pt]
  \item \href{../00-概念与术语/01-Agent术语表与概念边界.md}{Agent 术语表与概念边界}
  \item \href{../../../知识地图总索引.md}{知识地图总索引}
\end{itemize}

\end{document}
